\section*{Calibration of Merton and Kou Jump-Diffusion Models to QQQ}

\subsection*{Data and Methodology}

We calibrate the Merton (1976) and Kou (2002) jump-diffusion models to a
snapshot of QQQ option prices taken on 29~May~2026, with spot price
$S_0 = \$735.60$ and a risk-free rate of $r = 5.00\%$. QQQ ETF options
are American-style; we approximate them as European because the
calibration uses only out-of-the-money quotes and excludes deep
in-the-money contracts, where the early-exercise premium is largest.
This approximation is documented as a limitation in the closing
subsection. The raw option chain was retrieved through
\texttt{yfinance} for eight maturities spread across the term structure
(one week, two weeks, one month, two months, three months, seven
months, ten months, thirteen months); 2{,}480 quotes survived the
basic data-quality filters of \texttt{src/data\_pipeline.py} (positive
bid-or-last price, non-stale, and IV plausibly below 300\%).

For calibration we restrict to a clean near-ATM subset by applying four
additional filters: time to expiry $T \geq 0.05$ years (drops the two
shortest weekly maturities), out-of-the-money quotes only (puts with
$K < S_0$ and calls with $K > S_0$), a moneyness band $K/S_0 \in
[0.85, 1.15]$, and a Brent-recomputed implied volatility in $[0.10,
0.60]$. After cleaning, $M = 424$ quotes across six maturities
remained, ranging from $T = 0.077$~y (one month) to $T = 1.052$~y
(thirteen months).

For each candidate parameter vector $\theta$, model prices are computed
via the COS method of Fang and Oosterlee (2008) with
$N_{\text{COS}} = 512$ expansion terms and truncation parameter
$L = 10$. Each model price is inverted back to its Black--Scholes
implied volatility via Brent's method, and we minimise the vol-space
root mean squared error
\[
   \text{RMSE}(\theta) \;=\;
   \sqrt{\frac{1}{M}\sum_{i=1}^{M}
            \bigl( \sigma^{\text{IV},\text{model}}_i(\theta)
                   - \sigma^{\text{IV},\text{market}}_i \bigr)^{\!2}},
\]
where $M$ is the number of cleaned quotes (here $M = 424$) and
$N_{\text{COS}}$ controls the COS Fourier-series accuracy.

Optimisation runs in two stages. A \texttt{differential\_evolution}
global search (population size~12, 40 generations, evaluated on a
stratified sub-sample of 96 quotes for speed) locates the right region
in parameter space; an \texttt{L-BFGS-B} local refinement then polishes
the optimum on the full 424 quotes. All parameter bounds and optimiser
settings are read from \texttt{.env} through the
\texttt{CalibrationConfig} dataclass.

We acknowledge two limitations on the data side. The snapshot was taken
outside NYSE trading hours, so 98\% of cleaned quotes use last-trade
price as a fallback for mid-price. We mitigate this by recomputing all
implied volatilities ourselves via Brent rather than trusting
vendor-reported IVs, which are unreliable for stale or zero-bid quotes.
QQQ also pays a quarterly dividend (\textasciitilde$0.6\%$ annualised);
we do not subtract a dividend yield $q$ from the discount rate. This
biases the implied $\sigma$ marginally upward for longer-dated options;
a production calibration should be performed on forward prices.


\subsection*{Merton Calibration}

The calibrated Merton parameters are summarised in the table below.
The achieved vol-space RMSE is 0.01400, an average pricing error of
1.40 percentage points in implied volatility across the 424-quote
surface.

\begin{table}[H]
\centering
\caption{Calibrated Merton parameters (QQQ, 2026-05-29).}
\begin{tabular}{lrl}
\toprule
Parameter      & Value     & Interpretation \\
\midrule
$\sigma$       & $0.1858$  & continuous diffusion volatility \\
$\lambda$      & $0.1633$  & risk-neutral jump intensity (jumps per year) \\
$\mu_J$        & $-0.4954$ & mean log-jump size (at lower bound $-0.50$) \\
$\sigma_J$     & $0.3339$  & log-jump size standard deviation \\
$\kappa$       & $-0.3558$ & mean percentage jump $\mathbb{E}[e^J - 1]$ \\
RMSE           & $0.01400$ & vol-space RMSE on 424 quotes \\
\bottomrule
\end{tabular}
\end{table}

Two features deserve immediate comment. First, the diffusion volatility
$\sigma \approx 19\%$ is consistent with the realised volatility of QQQ
over the past three years and with the at-the-money level of the IV
surface. Indeed, this calibrated $\sigma$ is essentially identical to
the annualised realised volatility of QQQ log-returns over the
three-year W5 sample window ($\sqrt{252}\cdot 0.01242 \approx 19.7\%$),
providing an independent cross-check that the calibration recovers the
diffusion regime correctly. Second, the optimum lands on the lower bound of the project's
specified $\mu_J$ range: even when the bound is widened to $-0.50$, the
optimiser still attaches itself to it, implying a calibrated model in
which crashes are very rare ($\lambda = 0.16$, i.e., roughly one event
every six years) but catastrophic in magnitude ($\kappa \approx -36\%$
on average). This boundary behaviour is the empirical signature of the
identification problem analysed in the next subsection.

\safefig{../outputs/w6/figures/smile_fit_shortest.png}{Market vs Merton
and Kou implied volatilities for the one-month expiry. Black dots:
market quotes; solid line: Merton fit; dashed line: Kou
fit.}{fig:w6-smile-short}

\safefig{../outputs/w6/figures/smile_fit_all_maturities.png}{Smile fit
across all six calibrated maturities.}{fig:w6-smile-all}


\subsection*{The Identification Problem}

The project specification asks us to demonstrate that $\sigma$ and
$\lambda$ are not jointly identified from option prices. Our calibration
produces a stronger result: the more severe non-identification is
between $\lambda$ and $\mu_J$, evidenced by the boundary behaviour above
and confirmed by three independent diagnostics.

\subsubsection*{Diagnostic 1: Boundary Attachment}

The unregularised optimum lies at $\mu_J = -0.4954$, i.e., effectively
on the lower bound $\mu_J = -0.50$. We did not impose this bound
\emph{a priori} to constrain the problem; it was the widest economically
plausible range. Nevertheless, the optimiser is unwilling to move away
from it, indicating that the data prefers a corner solution.

\subsubsection*{Diagnostic 2: Two-Dimensional Loss Surface}

Fixing $(\sigma, \sigma_J)$ at their calibrated values and sweeping an
$18 \times 18$ grid in $(\lambda, \mu_J)$ on the sub-sampled dataset
reveals a curved low-loss valley running from the bottom-left to the
upper-right of the plane. Within RMSE $+\,0.002$ of the grid minimum we
find 7 distinct parameter pairs, all producing economically
indistinguishable fits. The market data does not contain enough
information to discriminate among them.

\safefig{../outputs/w6/loss_surface_lambda_muJ.png}{Merton vol-space RMSE
surface in $(\log_{10}\lambda, \mu_J)$ with $(\sigma, \sigma_J)$ fixed
at calibrated values. White contours show iso-loss lines; the red star
marks the calibrated optimum.}{fig:w6-loss-surface}

\subsubsection*{Diagnostic 3: L1-Regularisation Path}

Adding a penalty $\alpha\,|\mu_J|$ to the loss pulls the optimum away
from the boundary, as expected for an ill-posed inverse problem. The
regularisation path is computed on the stratified 96-quote sub-sample
used for the global search to keep its five calibrations tractable in
wall-clock time; the $\alpha = 0$ row is therefore numerically close
to, but not exactly identical to, the full-surface main optimum in
Table above. The qualitative trend (boundary attachment and recovery
under penalty) is unaffected.

\begin{table}[H]
\centering
\caption{Regularisation path. The L1 penalty $\alpha\,|\mu_J|$ added to
the loss pushes the optimum out of the boundary; as $\alpha$ grows the
optimum approaches the literature-consistent regime at a modest cost
in fit.}
\begin{tabular}{rrrrrr}
\toprule
$\alpha$ & $\sigma$ & $\lambda$ & $\mu_J$   & $\sigma_J$ & RMSE (unreg.) \\
\midrule
$0.000$  & $0.1863$ & $0.163$   & $-0.5000$ & $0.3051$   & $0.01398$ \\
$0.005$  & $0.1794$ & $0.328$   & $-0.2718$ & $0.2493$   & $0.01448$ \\
$0.010$  & $0.1741$ & $0.486$   & $-0.1967$ & $0.2161$   & $0.01502$ \\
$0.020$  & $0.1659$ & $0.803$   & $-0.1329$ & $0.1791$   & $0.01593$ \\
$0.050$  & $0.1369$ & $2.502$   & $-0.0606$ & $0.1164$   & $0.01851$ \\
\bottomrule
\end{tabular}
\end{table}

At $\alpha = 0.05$ the solution lands at $\mu_J \approx -0.06$ and
$\lambda \approx 2.5$, the literature-consistent regime of ``moderate
jumps several times per year'' documented in Andersen, Benzoni and Lund
(2002) and Eraker, Johannes and Polson (2003). The cost is a 32\%
increase in RMSE ($0.01400 \to 0.01851$), which we view as a favourable
trade-off: a modest deterioration in fit buys parameters that are both
economically plausible and statistically stable.

\safefig{../outputs/w6/regularisation_path.png}{Regularisation path. Each
panel shows one Merton parameter as a function of the L1 penalty
strength $\alpha$. Bottom-right: the unregularised RMSE deteriorates
only modestly as $\alpha$ grows.}{fig:w6-reg-path}

The economic content of this exercise is the \emph{peso problem} of
Backus, Chernov and Martin (2011): the smile of an equity index can be
generated either by frequent, mild jumps or by rare, catastrophic ones,
and the data alone cannot tell the two regimes apart.


\subsection*{Kou Comparison}

We re-calibrated the same QQQ surface to the Kou (2002)
double-exponential jump model, which replaces Merton's normal log-jump
with the asymmetric density
\[
   f_J(j) \;=\; p\,\eta_1 e^{-\eta_1 j}\,\mathbf{1}_{j \geq 0}
              + q\,\eta_2 e^{\eta_2 j}\,\mathbf{1}_{j < 0},
   \qquad p + q = 1,\ \eta_1 > 1,\ \eta_2 > 0.
\]
The characteristic function from W2 substitutes directly into the COS
kernel; calibration uses the same two-stage protocol.

\begin{table}[H]
\centering
\caption{Calibrated Kou parameters (QQQ, 2026-05-29). The
``$\approx 41\%$'' figure for $1/\eta_2$ refers to the mean
\emph{log}-jump magnitude on the downside; the corresponding expected
percentage price drop conditional on a downward jump is
$1 - \eta_2/(\eta_2 + 1) \approx 29\%$.}
\begin{tabular}{lrl}
\toprule
Parameter      & Value     & Interpretation \\
\midrule
$\sigma$       & $0.1813$  & diffusion volatility \\
$\lambda$      & $1.825$   & jump intensity (jumps per year) \\
$p$            & $0.875$   & probability of upward jump ($q = 0.125$) \\
$\eta_1$       & $46.17$   & upward tail; mean upward log-jump $1/\eta_1 \approx 2.2\%$ \\
$\eta_2$       & $2.44$    & downward tail; mean downward log-jump magnitude $1/\eta_2 \approx 41\%$ \\
RMSE           & $0.01398$ & vol-space RMSE on 424 quotes \\
\bottomrule
\end{tabular}
\end{table}

The economic structure is striking: 87.5\% of jumps are small upward
moves (mean log-jump $\approx 2.2\%$), but the remaining 12.5\% are
large downward moves whose mean log-magnitude is about $41\%$ --- which
corresponds to an expected percentage drop, conditional on a down-jump,
of $1 - \eta_2/(\eta_2 + 1) \approx 29\%$. The Kou compensator is
\[
   \kappa^Q \;=\; \frac{p\,\eta_1}{\eta_1 - 1}
                 + \frac{q\,\eta_2}{\eta_2 + 1} - 1
              \;=\; \frac{0.875 \cdot 46.17}{45.17}
                  + \frac{0.125 \cdot 2.44}{3.44} - 1
              \;\approx\; -1.7\%.
\]
The mean percentage jump is small in magnitude because the frequent
small upward jumps and the rare large downward jumps approximately
offset each other in mean. The variance contribution, by contrast, is
dominated by the heavy left tail.

\subsubsection*{Goodness-of-Fit Comparison}

The table below places Merton and Kou side by side. Kou's vol-space
RMSE is marginally lower (0.01398 vs 0.01400), but the Akaike
information criterion marginally favours Merton ($\Delta\text{AIC} =
1.06$, with AIC $= -3612.02$ vs $-3610.96$); since $\Delta\text{AIC} < 2$,
the two models should be regarded as statistically very close rather
than decisively separated by this criterion. The peso-style
identification problem discussed below offers a more economically
meaningful interpretation of the comparison.

\begin{table}[H]
\centering
\caption{Model comparison. AIC penalises Kou for the extra parameter
despite a marginally better in-sample fit.}
\begin{tabular}{lrrr}
\toprule
Model  & $k$ params & RMSE     & AIC        \\
\midrule
Merton & $4$        & $0.01400$ & $-3612.02$ \\
Kou    & $5$        & $0.01398$ & $-3610.96$ \\
\bottomrule
\end{tabular}
\end{table}

This result is not the textbook expectation that Kou should dominate
Merton on equity surfaces; instead it reflects identification on the
model-class level. Merton's $\mu_J \approx -0.50$,
$\lambda \approx 0.16$ optimum (rare large jumps) and Kou's
$1/\eta_2 \approx 0.41$, $\lambda \approx 1.83$ optimum (occasional
very large downward jumps) are two different mathematical languages
describing the same underlying realisations. The QQQ surface, taken on
a single day, does not contain enough information to discriminate the
two regimes.

\safefig{../outputs/w6/figures/rmse_by_maturity.png}{Per-maturity RMSE:
Merton (blue) vs Kou (orange).}{fig:w6-rmse-maturity}


\subsection*{Jump Risk Premium}

The final element of the calibration compares the risk-neutral jump
parameters from the option calibration with the physical parameters
from the W5 jump-detection pipeline. The W5 BNS-style diagnostic
produces two jump-day classifications: a \emph{liberal} flag (Z-score
above 2.0, 99 jump days) and a \emph{conservative} flag (Z-score above
2.0 \emph{and} return above $3\times$ rolling-vol, only 4 jump days).
We report the premium against both, since neither corresponds exactly
to the events the risk-neutral $\lambda^Q$ is pricing.

\subsubsection*{Intensity}

The risk-neutral intensity is $\lambda^Q = 0.163$ jumps per year. The
liberal physical intensity is $\lambda^{P,\text{lib}} = 34.08$ and the
conservative one $\lambda^{P,\text{cons}} = 1.38$. A direct difference
$\lambda^Q - \lambda^P$ is meaningless because the three quantities
count different events: liberal-$P$ counts any volatility-test
rejection (average $|r| \approx 1.3\%$), conservative-$P$ counts only
the very largest moves (average $|r| \approx 6.3\%$), while
$\lambda^Q$ measures the intensity of much rarer events
($|\kappa^Q| \approx 36\%$) required to bend the IV smile.

\subsubsection*{Jump-Variance Risk Premium}

The cleaner, units-free metric is the ratio of total annualised jump
variance under each measure:
\[
   \text{JVRP} \;=\;
   \frac{\lambda^Q \cdot \mathbb{E}^Q[J^2]}
        {\lambda^P \cdot \mathbb{E}^P[J^2]}.
\]
Plugging in $\mathbb{E}^Q[J^2] = \mu_J^2 + \sigma_J^2 = 0.357$, and the
empirical second moments $\mathbb{E}^{P,\text{lib}}[J^2] =
3.7 \times 10^{-4}$, $\mathbb{E}^{P,\text{cons}}[J^2] = 3.96 \times
10^{-3}$:
\[
   \text{JVRP}_{\text{lib}}  \;=\;
      \frac{0.163 \times 0.357}{34.08 \times 3.7 \times 10^{-4}}
      \;=\; \frac{0.0583}{0.0127} \;=\; 4.57,
\]
\[
   \text{JVRP}_{\text{cons}} \;=\;
      \frac{0.163 \times 0.357}{1.38 \times 3.96 \times 10^{-3}}
      \;=\; \frac{0.0583}{0.0055} \;=\; 10.69.
\]

Investors price jump variance at between 4.6 and 10.7 times its
historical realisation, depending on how broadly ``a jump'' is defined.
Both numbers are large positive premia and consistent in sign with the
empirical estimates of Pan (2002), Eraker (2004), and Bollerslev and
Todorov (2011) for the S\&P~500, although our magnitudes are at the
upper end of the published range, a direct consequence of the
peso-style calibration with very rare but very large risk-neutral
jumps.

This finding directly underwrites the W8 alpha strategy:
short-volatility positions in QQQ index options harvest, on average,
the jump variance premium documented here, while bearing the
corresponding jump risk in the left tail of their P\&L distribution.


\subsection*{Limitations and W7 Hand-Off}

Our calibration has five limitations worth flagging for the volatility
analysis in W7 and the hedging study in W8.

\begin{enumerate}[leftmargin=*]
   \item \textbf{American-style options approximated as European.}
   QQQ ETF options are American. We use European COS pricing because
   the calibration restricts to out-of-the-money quotes, where the
   early-exercise premium is small. A production implementation
   should add an American-style correction (Bjerksund--Stensland) or
   work on forwards.

   \item \textbf{No dividend yield.} The QQQ ETF pays a small
   quarterly dividend (\textasciitilde$0.6\%$ annualised) that we do
   not subtract from the discount rate. This biases the implied
   $\sigma$ marginally upward for longer-dated options. The proper
   fix is to calibrate on the forward $F = S_0 e^{(r-q)T}$ rather
   than the spot.

   \item \textbf{Boundary behaviour and identification.} The Merton
   optimum lies on the lower bound for $\mu_J$ and would walk further
   if allowed. All downstream Greeks, hedging ratios, and P\&L
   distributions should be reported at \emph{both} the unregularised
   parameters and at the regularised parameters of the previous
   subsection, to avoid relying on a single corner solution.

   \item \textbf{Short-maturity smile.} Both Merton and Kou fit the
   shortest maturity systematically worse than longer-dated options.
   W7's smile-decay analysis should document this and point to richer
   models (Bates, CGMY) that address it.

   \item \textbf{Single-day snapshot outside trading hours.} Only 2\%
   of our quotes use a simultaneous live mid; the remainder fall back
   to last-trade prices. A multi-date study during NYSE hours would
   tighten all conclusions, particularly the Kou-vs-Merton comparison.
\end{enumerate}

The calibrated parameters are exposed to W7 through
\texttt{src/iv\_surface.py}, which provides:
\begin{itemize}[leftmargin=*]
   \item \texttt{get\_merton\_params()} and \texttt{get\_kou\_params()},
         returning the calibrated dictionaries;
   \item \texttt{get\_merton\_params\_regularised(alpha)}, returning the
         regularised alternatives at any $\alpha$ from the path above;
   \item \texttt{merton\_iv(...)}, \texttt{kou\_iv(...)}, returning the
         model implied volatility of a single option;
   \item \texttt{merton\_iv\_surface(S0, r, K\_grid, T\_grid, params)},
         returning a two-dimensional model IV grid suitable for direct
         comparison with the market surface in W7.
\end{itemize}
